\cotefile{151}
\scanpage{32}
de \uline{locale trivialité normale} a un rôle
à jouer (celle de lissité \unsure{restant} comme
devant techniquement inutile \deleted{\unsure{?}}, son rôle
était plutôt de nature heuristique ...).

\deleted{\unsure{On}} On doit se débrouiller, dans ce cas,
pour démontrer que les propriétés homo-
topiques de l'inclusion $V^*_{Y,X}\subset V_{Y,X}$
\margin{\diagram{\begin{tikzcd}[row sep=small]
{} \\
Y \arrow[u, "{\unsure{\text{homo-topisme}}}"]
\end{tikzcd}}\deleted{\unsure{Id}}}
sont celles d'une \uline{fibration}
\uline{localement triviale de fibre} $V^*_{z,Z}$ :

\noindent (11) $V^*_{Y,X}\longrightarrow V_{Y,X}$ est une ``fibration
homotopique'' de fibre $V^*_{z,X}$,
de base $Y\xrightarrow{\approx\,\mathrm{homot.}} V_{Y,X}$

--- et c'est cela qui devrait même,
dans le contexte toposique (p. ex. celui
des schémas avec la top. étale) servir
de \uline{définition} de la ``locale trivialité \inserted{normale}
homotopique'' de l'inclusion $Y\subset X$.
(Bien sûr, dans le contexte schématique,
il faudra de plus travailler avec
des types d'homotopie profinis, et
même \inserted{sans doute} ``localiser'' ces types d'homotopie
\mspage{31}
